\chapter[Seguridad en PostgreSQL]{\emj{🔐}~Seguridad en PostgreSQL}

\section[Row-Level Security (RLS)]{\emj{🛡️}~Row-Level Security (RLS)}

RLS permite que la BD misma filtre filas según el usuario actual, sin que la aplicación tenga que implementar los filtros.
Ideal para arquitecturas multi-tenant donde cada tenant solo debe ver sus datos.

\begin{teoriabox}{Cómo funciona RLS}
\begin{itemize}
  \item Habilitado por tabla con \texttt{ALTER TABLE ... ENABLE ROW LEVEL SECURITY}
  \item Las policies definen qué filas puede ver/modificar cada usuario o rol
  \item Un superuser o el owner de la tabla \textbf{no está sujeto a RLS} por defecto (a menos que se use \texttt{FORCE ROW LEVEL SECURITY})
  \item Las policies se expresan como expresiones SQL que se agregan automáticamente a los WHERE
\end{itemize}
\end{teoriabox}

Sin RLS, el aislamiento multi-tenant depende completamente de que la aplicación filtre correctamente por \texttt{tenant\_id} en cada query.
Un bug, un olvidar agregar el filtro, o un path de código poco usado puede exponer datos de un tenant a otro.
Con RLS, el filtro vive en la base de datos y es imposible saltárselo desde la aplicación: incluso si alguien ejecuta \texttt{SELECT * FROM orders} directamente desde psql con las credenciales de la app, solo verá las filas de su tenant.

La aplicación debe establecer el tenant ID al inicio de cada transacción usando \texttt{SET LOCAL}, que garantiza que la variable se resetea al hacer commit o rollback, sin posibilidad de que un tenant ``quede pegado'' en la sesión.

\begin{lstlisting}[caption={RLS para multi-tenant SaaS}]
-- Habilitar RLS en la tabla:
ALTER TABLE orders ENABLE ROW LEVEL SECURITY;

-- Policy de SELECT: cada usuario solo ve sus propias órdenes
CREATE POLICY orders_tenant_isolation ON orders
  FOR SELECT
  USING (tenant_id = current_setting('app.current_tenant_id')::BIGINT);

-- Policy de INSERT: solo puede insertar para su propio tenant
CREATE POLICY orders_tenant_insert ON orders
  FOR INSERT
  WITH CHECK (tenant_id = current_setting('app.current_tenant_id')::BIGINT);

-- Policy para UPDATE y DELETE:
CREATE POLICY orders_tenant_update ON orders
  FOR UPDATE
  USING (tenant_id = current_setting('app.current_tenant_id')::BIGINT);

-- La app debe setear el tenant ID al inicio de cada sesión/transacción:
SET app.current_tenant_id = '42';
-- O por transacción (más seguro):
BEGIN;
  SET LOCAL app.current_tenant_id = '42';
  SELECT * FROM orders;  -- automáticamente filtrado a tenant 42
COMMIT;

-- Bypassar RLS para el admin (superuser ya bypasea por defecto):
ALTER TABLE orders FORCE ROW LEVEL SECURITY;  -- aplica incluso al owner
-- Crear rol admin que bypasea:
CREATE ROLE admin_role BYPASSRLS;
\end{lstlisting}

\begin{lstlisting}[caption={Verificar policies activas}]
SELECT schemaname, tablename, policyname, permissive, roles, cmd, qual
FROM pg_policies
WHERE tablename = 'orders';

-- Testear como usuario específico:
SET ROLE app_user;
SET app.current_tenant_id = '42';
SELECT COUNT(*) FROM orders;  -- debe mostrar solo órdenes de tenant 42
RESET ROLE;
\end{lstlisting}

\section[Cifrado con pgcrypto]{\emj{🔑}~Cifrado con pgcrypto}

\texttt{pgcrypto} es la extensión de PostgreSQL para operaciones criptográficas dentro de la base de datos.
Permite dos categorías de uso:

\begin{itemize}
  \item \textbf{Hashing de contraseñas:} usando bcrypt (\texttt{crypt} + \texttt{gen\_salt('bf')}). El hash es unidireccional e incluye el salt, por lo que no hay que gestionar el salt por separado. La verificación se hace comparando el hash almacenado con el hash del intento usando la misma función.
  \item \textbf{Cifrado simétrico de datos sensibles:} usando AES via \texttt{pgp\_sym\_encrypt}. A diferencia del hashing, el cifrado es reversible con la clave. Útil para datos que la aplicación necesita leer en claro (números de tarjeta, datos médicos) pero que no deben ser legibles si alguien accede directamente al archivo de datos o al backup.
\end{itemize}

La clave de cifrado \textbf{nunca} debe aparecer en el código o en la configuración del repositorio.
Debe provenir de un gestor de secretos (HashiCorp Vault, AWS Secrets Manager, GCP Secret Manager) inyectada en la aplicación en tiempo de ejecución.

\begin{lstlisting}[caption={Cifrado simétrico y hashing con pgcrypto}]
CREATE EXTENSION pgcrypto;

-- Hashing de contraseñas (bcrypt):
-- Guardar hash:
INSERT INTO users (email, password_hash)
VALUES ('user@example.com', crypt('mi_password_123', gen_salt('bf', 12)));
-- 'bf' = blowfish (bcrypt), 12 = cost factor

-- Verificar contraseña:
SELECT id FROM users
WHERE email = 'user@example.com'
  AND password_hash = crypt('mi_password_123', password_hash);
-- Si retorna fila → contraseña correcta

-- Cifrado simétrico de datos sensibles (AES):
-- Cifrar al insertar:
INSERT INTO payments (id, card_number_encrypted)
VALUES (1, pgp_sym_encrypt('4111111111111111', 'clave_de_cifrado_secreta'));

-- Descifrar:
SELECT pgp_sym_decrypt(card_number_encrypted, 'clave_de_cifrado_secreta')
FROM payments WHERE id = 1;

-- MEJOR PRÁCTICA: no guardar la clave en la query. Pasar desde app.
-- La clave debe venir del gestor de secretos (Vault, AWS Secrets Manager)
-- y nunca aparecer en los logs de queries

-- Cifrado asimétrico (PGP):
-- Generar claves con gpg externo, usar pgp_pub_encrypt / pgp_pub_decrypt
\end{lstlisting}

\section[Autenticación y pg\_hba.conf]{\emj{🔒}~Autenticación y pg\_hba.conf}

\begin{lstlisting}[language=, caption={pg\_hba.conf — reglas de autenticación}]
# Formato: TYPE  DATABASE  USER  ADDRESS  METHOD
# Las reglas se evalúan de arriba hacia abajo; primera coincidencia aplica

# Rechazar conexiones locales de superuser por red:
hostssl  all  postgres  0.0.0.0/0  reject

# Requerir SSL + scram-sha-256 para conexiones de red:
hostssl  all  all  0.0.0.0/0  scram-sha-256

# Conexión local sin contraseña (para mantenimiento desde el servidor):
local  all  postgres  peer

# Replicación solo desde servidor de standby específico:
hostssl  replication  replicator  10.0.0.2/32  scram-sha-256

# md5 está deprecado (vulnerable a MITM); usar scram-sha-256
\end{lstlisting}

\begin{lstlisting}[caption={SSL/TLS — requerir cifrado en tránsito}]
-- postgresql.conf:
ssl = on
ssl_cert_file = 'server.crt'
ssl_key_file = 'server.key'
ssl_ca_file = 'root.crt'     -- para verificación de clientes (mTLS)

-- Verificar que la conexión usa SSL:
SELECT ssl, version, cipher FROM pg_stat_ssl WHERE pid = pg_backend_pid();

-- Forzar SSL para un usuario:
ALTER USER app_user SET ssl = 'on';  -- no existe directamente; usar pg_hba.conf hostssl

-- En el cliente (psql):
-- psql "sslmode=verify-full sslrootcert=root.crt ..."
\end{lstlisting}

\section[Principio de Mínimo Privilegio]{\emj{👤}~Principio de Mínimo Privilegio}

El principio de mínimo privilegio (Principle of Least Privilege) establece que cada componente del sistema debe tener solo los permisos estrictamente necesarios para su función.
En el contexto de PostgreSQL, significa que el usuario con el que se conecta la aplicación no debe poder hacer DROP TABLE, ALTER TABLE, o borrar datos de otras bases.
Un atacante que compromete las credenciales de la aplicación solo puede hacer lo que la aplicación puede hacer: el daño queda acotado.

La estructura recomendada es usar \textbf{roles} como capas de permisos (no usuarios directamente), y asignar usuarios a roles.
Esto permite revocar y conceder permisos a múltiples usuarios de una vez, y mantener políticas consistentes.

\begin{lstlisting}[caption={Roles y privilegios granulares}]
-- Crear roles específicos por función:
CREATE ROLE app_readonly;
GRANT CONNECT ON DATABASE mydb TO app_readonly;
GRANT USAGE ON SCHEMA public TO app_readonly;
GRANT SELECT ON ALL TABLES IN SCHEMA public TO app_readonly;
ALTER DEFAULT PRIVILEGES IN SCHEMA public GRANT SELECT ON TABLES TO app_readonly;

CREATE ROLE app_writer;
GRANT app_readonly TO app_writer;
GRANT INSERT, UPDATE, DELETE ON ALL TABLES IN SCHEMA public TO app_writer;
ALTER DEFAULT PRIVILEGES IN SCHEMA public
  GRANT INSERT, UPDATE, DELETE ON TABLES TO app_writer;

-- Usuario de la aplicación hereda del rol:
CREATE USER app_user WITH PASSWORD 'securepass';
GRANT app_writer TO app_user;

-- Nunca conectar la app como superuser o como el owner del schema
-- El owner puede DROP TABLE, ALTER TABLE, etc.

-- Revocar privilegios no necesarios:
REVOKE CREATE ON SCHEMA public FROM PUBLIC;  -- PG 14+ lo hace por defecto
REVOKE ALL ON DATABASE mydb FROM PUBLIC;
\end{lstlisting}

\section[pg\_audit: logging de queries]{\emj{📋}~pg\_audit: logging de queries}

Los logs estándar de PostgreSQL registran conexiones y queries lentas, pero no tienen suficiente detalle para auditorías de compliance (PCI-DSS, SOC 2, HIPAA).
\texttt{pgaudit} es una extensión que añade logging granular de todas las operaciones: qué tablas se leyeron, qué datos se modificaron, qué DDL se ejecutó, y qué usuario lo hizo.
Los logs de auditoría deben enviarse a un sistema externo (SIEM, S3 con cifrado) donde la BD misma no pueda modificarlos, garantizando su integridad para auditorías.

\begin{lstlisting}[language=, caption={Configurar pg\_audit para compliance}]
# postgresql.conf:
shared_preload_libraries = 'pgaudit'
pgaudit.log = 'write, ddl'   # ddl: CREATE/ALTER/DROP; write: INSERT/UPDATE/DELETE
pgaudit.log_catalog = off    # no loggear queries de sistema (pg_catalog)
pgaudit.log_relation = on    # incluir nombre de tabla/vista
pgaudit.log_parameter = on   # incluir parámetros de prepared statements

# Por rol específico:
# ALTER ROLE audited_user SET pgaudit.log = 'all';

# Formato del log:
# AUDIT: SESSION,1,1,DDL,CREATE TABLE,TABLE,public.users,...
# AUDIT: SESSION,2,1,WRITE,INSERT,TABLE,public.orders,...
\end{lstlisting}

\begin{entrevistabox}
\begin{enumerate}
  \item ¿Qué es Row-Level Security y cuándo la usarías?
  \item Explica la diferencia entre \texttt{USING} y \texttt{WITH CHECK} en una policy de RLS.
  \item ¿Por qué md5 está deprecado en pg\_hba.conf? ¿Qué usar en su lugar?
  \item Un desarrollador quiere conectar la app a PG como superuser ``por simplicidad''. ¿Qué le dices?
  \item ¿Cómo almacenarías contraseñas de usuarios en PostgreSQL? ¿Y números de tarjeta de crédito?
  \item ¿Qué es mTLS y cuándo lo usarías en una conexión a PostgreSQL?
\end{enumerate}
\end{entrevistabox}
